\uuid{D3xz}
\titre{Logique}
\auteur{bodin ; croizat ; sacre}
\organisation{exo7}
\licence{CC BY-NC-SA 4.0}
\datecreate{2022-03-08}
\niveau{L1}
\module{Algèbre}
\chapitre{Logique, ensemble, raisonnement}
\sousChapitre{Logique}
\difficulte{3}
\theme{Logique}
\exo7qcmid{hilisit.logique.014}

\contenu{
\begin{question}
\qclassification{100.01}
Soit $x\in\Rr$.
Pour quelles phrases, l'assertion est vraie si on remplace "$??$" par "$\exists$", mais est fausse si on remplace "$??$ " par "$\forall$" ?
\begin{answers}
	\good{$??\; x\in \Rr \quad x^2>0$}
	\bad{$??\; x\in \Rr \quad x^2-2x+1\ge 0$}
	\good{$??\; x\in\Rr \quad x^2-2x+1=0$}
	\good{$??\; x\in\Rr \quad x^2\le 0$}
\end{answers}
\begin{explanations}
Il existe un réel $x$ tel que $x^2>0$, par exemple $x=1$, mais ce n'est pas vrai pour $x=0$.\\
$x^2-2x+1=(x-1)^2\ge 0$ pour tout réel $x$.\\
 $x^2-2x+1$ est nul pour $x=1$ mais n'est pas nul pour $x=0$. \\
 $x^2\le 0$ pour $x=0$ mais ce n'est pas vrai pour $x=1$.
\end{explanations}
\end{question}
}
